\documentclass[11pt]{article}
\usepackage[T1]{fontenc}
\usepackage{textcomp}
\usepackage[margin=0.75in]{geometry}
\usepackage{amsmath,amssymb,amsthm}
\usepackage{array,booktabs}
\usepackage{hyperref}
\setlength{\parindent}{0pt}
\newtheorem{theorem}{Theorem}
\theoremstyle{definition}
\newtheorem{definition}{Definition}
\title{Flow Proof Artifact: List-rev-cons-double}
\author{Generated by \texttt{flow doc proof}}
\date{Flow Proof Book}
\begin{document}
\maketitle
Reversing a double cons list.\\[0.5em]
\textbf{Source.} church — https://en.wikipedia.org/wiki/Reverse\_(list)\\[0.5em]
\bigskip
\noindent{\large \textbf{Derived fact 1.} \textit{rev(cons(a, cons(b, nil))) = cons(b, cons(a, nil))} \hfill List \cdot reverse \cdot reverse of double cons}\par
\smallskip\noindent\textit{Coordinate: List \cdot reverse \cdot reverse of double cons}\par
\smallskip\noindent\textit{Source: church}\par
\smallskip\noindent\textit{Built on: reverse of cons is append of reverse tail, for reverse on List}\par
\medskip
\noindent\textbf{Goal.} rev(cons(a, cons(b, nil))) = cons(b, cons(a, nil))\par
\noindent$\displaystyle \forall a \in \mathbb{N} \forall b \in \mathbb{N}\quad rev(cons(a, cons(b, nil))) = cons(b, cons(a, nil))$\par
\medskip
\noindent
\begin{tabular}{@{} >{\raggedright\arraybackslash}p{0.06\textwidth} >{\raggedright\arraybackslash}p{0.47\textwidth} >{\raggedleft\arraybackslash}p{0.41\textwidth} @{}}
\textbf{\#} & \textbf{Proof} & \textbf{Mathematics} \\
\midrule
\textbf{1.} & We prove that reverse of double cons for reverse on List. & \\[0.45em]
\textbf{2.} & We invoke the derived fact governing reverse on List: reverse of cons is append of reverse tail, for reverse on List (instantiated for a, cons(b, nil)). & \\[0.45em]
\textbf{3.} & From step 2, this implies rev(cons(a, cons(b, nil))) equals cons(b, cons(a, nil)). Hence proven. & $\displaystyle rev(cons(a, cons(b, nil))) = cons(b, cons(a, nil))$ \\[0.45em]
\bottomrule
\end{tabular}
\label{thm:List-reverse-reverse_of_double_cons}
\par\medskip\hrule
\end{document}
